\section{Communication and parallel decomposition}
\label{sec:C07:parallelism}
Data parallelism replicates a model across examples and combines gradients. Tensor parallelism partitions individual matrix operations. Pipeline parallelism partitions layers across stages. Expert parallelism distributes conditional experts. Sequence or context parallelism partitions token-related work and associated state. These decompositions can be combined, but their communication patterns interact with the physical topology. An arrangement that minimizes memory can perform poorly when frequent communication crosses a slow link. The parallelism configuration is part of the scientific method because it changes feasible batch shapes, numerical reduction order, and time available for training.

For a ring all-reduce of a tensor with $D_g$ bytes over $n$ ranks, an ideal algorithm transfers approximately $2(n-1)D_g/n$ bytes per rank, ignoring protocol overhead. With a hypothetical 14-GB gradient tensor and eight ranks, this is 24.5 GB per rank per all-reduce. The formula explains why gradient accumulation can amortize communication over several microbatches, provided the accumulation semantics are correct. It does not predict duration without topology, effective bandwidth, overlap, and contention. Multiple parallel groups sharing the same links can invalidate an isolated collective microbenchmark as a training-time predictor.

Megatron's tensor-parallel construction partitions complementary matrix dimensions so that communication can be placed between substantial blocks of computation~\cite{shoeybi2019megatron}. The pinned \path{megatron/core/tensor_parallel/layers.py} exposes \texttt{ColumnParallelLinear}, \texttt{RowParallelLinear}, and \texttt{VocabParallelEmbedding}~\cite{repo_megatron}. For $Y=XW$ with $W$ split by output columns, each rank computes a subset of output features. A following compatible row partition can consume those features locally before a reduction combines partial outputs. The exact collective placement depends on the chosen layer composition and implementation options; a generic claim that tensor parallelism simply divides all compute by the rank count is false.

Pipeline schedules trade utilization against activation residency and stale work. A simple fill-and-drain pipeline has bubbles whose relative impact decreases with more microbatches, while more microbatches can increase activation memory or delay updates. Interleaving and other schedules change the tradeoff. For long agent traces, variable sequence lengths complicate balance: a stage can wait on an unusually expensive attention computation even if each stage has the same layer count. Measure stage-level time distributions and communication stalls rather than inferring balance from parameter allocation alone.

